
\subsection{Le modèle des facteurs de vue d'ordre 1}
Les facteurs de vue d'ordre 1 repose sur l'hypothèse d'\textbf{uniformité de la radiosité et de la température sur une surface}. Cette hypothèse permet d'introduire des $\overline{J_i^1} \; [W/m^2]$, et de découpler le terme géométrique $f_{ij}$ de la radiosité dans la formule \ref{def:irra}.  On a alors pour $\forall i$:
\red{partir de l'équation précédente}
\begin{equation}
    G_i(\vec{r_i}) = \sum_{j = 1}^{n} \overline{J_j}^1 f_{ij}(\vec{r_i}) \quad \text{avec } f_{ij}(\vec{r_i}) =\frac{1}{\pi} \int_{A_i} d^2(\vec{r_j})\bigg[\frac{cos(\theta_{ij})cos(\theta_{ji})}{||\vec{r_i}-\vec{r_j}||^2} \chi{(\vec{r_i},\vec{r_j})}\bigg] 
\end{equation}

L'intégration sur la surface $i$ permet enfin d'obtenir une approximation à l'ordre 1 la fraction de radiation quittant la surface $i$ intercepté par la surface $j$ appelée facteur de vue.

\begin{equation}
    \label{fv}
    F_{ij} = \frac{1}{A_i} \int_{A_i} \int_{A_j} d^2(\vec{r_i}) d^2(\vec{r_j})\bigg[\frac{cos(\theta_{ij})cos(\theta_{ji})}{\pi ||\vec{r_i}-\vec{r_j}||^2} \chi(\vec{r_i},\vec{r_j})\bigg] 
\end{equation}


Les hypothèses faites reviennent à considérer un problème d'éclairement entre différentes surface couplées. La connaissance des facteurs de vue permet alors de reformuler l'équation \ref{def:radiosite} pour exprimer les radiosités en fonction de la variable température : 

\begin{equation}
    \overline{J_i^1} = \epsilon_i \sigma \overline{T}^4_i  + \rho_i \sum_{j=1}^{n} \overline{J_j^1} F_{ij}
\end{equation}


On remarque par ailleurs que les facteurs de vue respectent des propriétés physiques fondamentales, à savoir :

\begin{equation}    
    \left\lbrace
    \begin{array}{r c c c  }
    
      \forall i ,\forall j, A_i F_{ij} & = & A_j F_{ji} & \text{(relation de réciprocité)}\\
      \forall i , \sum_{j = 1}^n F_{ij} & = & 1 & \text{(relation de fermeture)}\\ 
      \forall i ,\forall j, F_{ij} & \geq & 0 & \text{(positivité)}
    \end{array}\right.
\end{equation}
\label{def:fv}

Enfin, l'Equation \ref{eq:qr} qui permet de définir le flux radiatif devient :
\begin{equation}
    \forall \vec{r}_i \in \text{Surface}_i,\quad q_i(\vec{r}_i) =  \overline{J}_i^1  -  \sum_{j=1}^{n} \overline{J_j^1} F_{ij} \quad = \quad \epsilon_i \sigma \overline{T}^4_i  - \epsilon_i \sum_{j=1}^{n} \overline{J_j^1} F_{ij}
    \label{eq:nice_qr}
\end{equation}

